% !TEX encoding = UTF-8 Unicode
% !TEX program = lualatexmk
% was lualatex
% don't forget Edit > Substitutions > Text Replacement
% needs file combining_preprocessor.lua, & BibTex.bib (e.g.) for \bibliography

\documentclass[12pt]{article}

% packages

\usepackage{amsmath}		% needed for align & subeqs; must load before unicode-math
\usepackage{bm}			% for bold math symbols

% LuaLaTeχ
\usepackage[math-style=ISO,colon=literal]{unicode-math}	% "italic" capital Greek, math colon = text colon
%	\setmathfont{XITS Math}					% has subscript kerning bug
	\setmathfont{STIX Two Math}					% best; Times-like; in TeXLive 2018
	\setmainfont{XITS}
%	\setmainfont{STIX Two Text}[FakeStretch=0.95]	%[OpticalSize=12]
	\setsansfont{TEX Gyre Heros}
%	\setmonofont{Consolas}
\AtBeginDocument{\directlua{require("combining_preprocessor.lua")}}	% combining characters

\usepackage{microtype}		% fewer hyphens, etc.

%\usepackage{graphicx}
\def\fig#1{\saveimageresource {#1}\useimageresource\lastsavedimageresourceindex}
\def\figscale#1#2{\saveimageresource width#2 {#1}\useimageresource\lastsavedimageresourceindex}

\def\rgboo#1{\pdfextension literal{#1 rg #1 RG}}
\def\rgbo#1#2{\rgboo{#1}#2\rgboo{0 0 0}}
\def\rgb#1#2{\mark{#1}\rgbo{#1}{#2}\mark{0 0 0}}

\usepackage{slashed}		% Feynman: see  ̸ below
\usepackage{cite}			% for hyphenated citations
\usepackage{hyperref}
\hypersetup{
	bookmarksnumbered,
	unicode,			% use with \texorpdfstring
	colorlinks,			% avoid stupid boxes
%	citecolor=[rgb]{.9,0,.5},	% \cite
	urlcolor=[rgb]{0,0,1},	% \href
%	linkcolor=[rgb]{0,.7,0}	% \ref , toc
	}

% fancy boxes for eqs
\usepackage{empheq}
\usepackage[most]{tcolorbox}					% messes up syncing for title page
\definecolor{realcyan}{rgb}{0,1,1}
\definecolor{realyellow}{rgb}{1,1,0}
\newtcbox{\mymath}[1][]{%
    nobeforeafter, math upper, tcbox raise base,
    enhanced, colframe=realcyan!50!white,
    colback=realyellow!5!white, boxrule=1pt, drop lifted shadow, sharp corners
    #1}

% text style parameters		textile?

\textheight=9in						% height of body
\textwidth=6.5in						% width of body
\topmargin=-.5in						% top margin (less 1")
\oddsidemargin=0in \evensidemargin=0in	% centered margins (less 1")
\def\baselinestretch{1.2}				% magnification for line spacing
\parskip=\medskipamount				% space between paragraphs

% text

%% to get equations numbered by subsection
%\catcode`@=11
%\@addtoreset{equation}{section}
%\catcode`@=12
%\numberwithin{equation}{subsection}
% to get equations numbered by section
\numberwithin{equation}{section}

\def\hang{\hangindent\parindent}
\def\bitem{\par\hang\textindent}
\def\bitemitem{\par\indent \hangindent2\parindent \textindent}
\def\textindent#1{\indent\llap{#1\enspace}\ignorespaces}

% math

\mathcode`\*="702A                  	% * now always complex conjugate
\def\f#1#2{{\textstyle{#1\over#2}}}	% fraction

% stop amsmath whining about \over
\catcode`@=11
\let\over\@@over
\catcode`@=12

% accent over:
\def\on#1#2{{\buildrel{{\mkern2.5mu\raise-.1em\hbox{$\scriptstyle#1$}\mkern-2.5mu}}\over{#2}}}
\def\ron#1#2{{\buildrel{{\raise-.1em\hbox{$\scriptstyle#1$}}}\over{#2}}}	% for unitalicized
\def\dt#1{\on{\hbox{\bf .}}{#1}}								% (big) dot over: see  ̇  below  
\def\ddt#1{\on{\hbox{\bf .\kern-1pt.}}#1}							% double dot: see ¨ below
\def\under#1#2{\mathop{\null#2}\limits_{#1}}						% accent under

% nachos
%\def\dd{\hbox{\,\Large$\triangleright$}}
%\def\dd{\hbox{\scriptsize$\triangleright$}}
\def\dd{\hbox{\large$\smalltriangleright$}}

% Die, gamma!
\def\dig#1{\setbox0=\hbox{$#1M$}
	\hskip.06\wd0 \vrule width.07\wd0 height.63\wd0 depth.01\wd0 
	\vrule width.37\wd0 height.63\wd0 depth-.56\wd0 \hskip-.4\wd0
	\vrule width.25\wd0 height.35\wd0 depth-.28\wd0 
	\vrule width.07\wd0 height.35\wd0 depth-.17\wd0 \hskip.14\wd0}
\def\digamma{{\mathpalette\dig{}}}

% otherwise undefined utf-8
%\catcode`∙=13	\def ∙{\cdot}		% Option-*
\catcode`‹=13	\def ‹{\langle}		% Option-[
\catcode`›=13	\def ›{\rangle}		% Option-]
\catcode`❴=13	\def ❴{\{}			% Option-{
\catcode`❵=13	\def ❵{\}}			% Option-}
\catcode`√=13	\def √{\sqrt}		% Option-M
\catcode`§=13	\def §{\section}		% Option-i
\catcode`¶=13	\def ¶{\subsection}	% Option-I
\catcode`▷=13	\def ▷{\dd}
\catcode`𐊥=13	\def 𐊥{\digamma}
% spaces
\catcode`⎵=13	\def ⎵{\nobreak\ }				% Option-0
\catcode`⸏=13	\def ⸏{\quad}					% Option-)
%\catcode`␣=13	\def ␣{\TextOrMath{\thinspace}{\,}}		% (open box)
\catcode`˽=13	\def ˽{\TextOrMath{\thinspace}{\,}}		% Option-space (modifier letter shelf)
\catcode`ī=13 \def ī{{\bar\imath}}	%to correct for i with dot and bar
%\catcode`j̄=13 \def j̄{{\bar\jmath}}	%to correct for j with dot and bar
\catcode`・=13 \def・{\bullet}

% combining characters
\def\̂#1{\TextOrMath{{#1̂}}{{\hat{#1}}}}		% Option-6
\def\̌#1{\TextOrMath{#1̌}{\check{#1}}}	% Option-V
\def\̃#1{\TextOrMath{#1̃}{\tilde{#1}}}	% Option-n
\def\́#1{\TextOrMath{#1́}{\acute{#1}}}	% Option-~
\def\̀#1{\TextOrMath{#1̀}{\grave{#1}}}	% Option-`
\def\̇#1{\TextOrMath{#1̇}{\dt{#1}}}		% Option-;
\def\̈#1{\TextOrMath{#1̈}{\ddt{#1}}}		% Option-:
\def\̆#1{\TextOrMath{#1̆}{\breve{#1}}}
\def\̄#1{\TextOrMath{#1̄}{\bar{#1}}}		% Option--
\def\̊#1{\TextOrMath{#1̊}{\mathring{#1}}}	% Option-A
\def\̧#1{\TextOrMath{#1̧}{\c{#1}}}	% Option-A
\def\⃗#1{\vec{#1}}					% Option-R
\def\⃡#1{\overleftrightarrow{#1}}		% Option-E
\def\⃖#1{\overleftarrow{#1}}			% Option-W
\def\̸#1{\slashed{#1}}				% Option-/
\def\̵#1{\slashed{#1}}				% Option-/
% with single character only
%\def\͞#1{\overline{#1}}				% (combining double macron)
\def\̅#1{\overline{#1}}				% Option-? (combining overline)
%\def\̱#1{\underline{#1}}			% (combining macron below)
\def\̲#1{\underline{#1}}				% Option-_ (combining lowline)
\def\̱#1{{\underline{#1}}}				% shift+option+h
%\def\͡#1{\widehat{#1}}				% (combining double inverted breve)
\def\᷍#1{\widehat{#1}}				% Option-^ (combining double circumflex above)
\def\͠#1{\widetilde{#1}}				% Option-N


% environments
\def\eq#1{\begin{equation}#1\end{equation}}
\def\al#1{\begin{align}#1\nonumber\end{align}}
\def\ga#1{\begin{gather}#1\nonumber\end{gather}}
\def\li#1{\begin{itemize}#1\end{itemize}}
\def\se#1{\begin{subequations}\begin{align}#1\end{align}\end{subequations}}
\def\ox#1{\begin{empheq}[box=\mymath]{equation}#1\end{empheq}}
\def\ma#1{\begin{pmatrix}#1\end{pmatrix}}
\def\ar#1{\begin{matrix}#1\end{matrix}}

\def\rm{\mathrm}	% but use as \rm{ }, not {\rm }
\def~{\sim}

%%%%%%%%%%%%%%%%%%%%%%%%%%

